\chapter{Representing a company's historical record}
\label{ch:history}

\section{The record is a collection of assertions}

The original record includes financial statements, notes, accounting policies,
auditor opinions, management discussion, segment tables, security terms,
regulatory returns and subsequent announcements. A complete ingestion process
accounts for every page and every extracted block. A complete financial model
is a stronger claim: it would require economically relevant information that
public reports may never disclose. AccDSL therefore reports document coverage,
semantic coverage and simulation coverage separately.

Each source span receives a disposition: structured observation, narrative
assertion, relationship, policy evidence, contextual material, duplicate, or
unresolved extraction. ``Unresolved'' remains visible. A commitment buried in a
note is not discarded because it has no current balance-sheet line. An auditor's
qualification is retained as evidence about the record; it does not become an
automatic numerical haircut. Management's forecast is a dated assertion with
an author and modality, not an observed future outcome.

An observation is an immutable record with the fields below. Values inferred
from observations are separate records, linked by a derivation. This avoids
overwriting a company's statement with an analyst's preferred reconstruction.

\begin{longtable}{L{35mm}L{101mm}}
\toprule Field group & Minimum contents\\\midrule\endhead
Identity and source & Observation ID; document hash; filing identifier;
PDF page and bounding box or HTML/XBRL locator; literal text; extraction version.\\
Reporting context & Legal entity; reported group perimeter; statement or note;
instant or interval; standard and taxonomy versions; dimensions such as segment,
product, geography, seniority or measurement category.\\
Amount & Original string; parsed number; sign interpretation; unit; scale;
currency; precision or rounding interval; missing-value category.\\
Meaning & Source label; taxonomy QName when present; one or more proposed
concept mappings; whether the assertion is reported, estimated or conditional.\\
Timing & Economic validity when known; reporting period; publication and
availability timestamps; ingestion timestamp; correction or supersession links.\\
Reliability and review & Extraction confidence, reviewer decision, source
status and known contradictions, all distinct from a probability over histories.\\
\bottomrule
\end{longtable}

A narrative assertion adds actor, predicate, object, scope, quantifier and
modality. For example, ``substantially all debt is fixed rate'' has a source
span and a qualitative quantifier. It does not assert $\alpha=1$. A reviewed
interpretation can introduce a numerical restriction such as
$\alpha\geq a$, but the threshold $a$ remains an analyst assumption unless
the source defines it. Boilerplate policy language also does not prove that a
company holds a particular instrument.

\section{Different names, different measurements}

An accounting concept is more than a row label. ``Debt'' can refer to legal
principal, amortised carrying amount, undiscounted contractual cash flows or a
fair-value disclosure. These amounts can all be correct at the same date.
Likewise, ``revenue'' may be reported gross or net, a segment measure may differ
from the consolidated measure, and an alternative performance measure may
exclude selected expenses. A mapping must state its relationship: equivalent
meaning, component of, broader than, presentation reclassification, or
measurement transformation.

The common vocabulary will preserve source concepts as first-class identifiers.
Company extensions may inherit a broad economic category without being declared
identical to a standardised concept. For instance, an extension can be a
component of borrowings while retaining an unresolved split between lease
liabilities and other financing. Many-to-many mappings are allowed; each
mapping has applicable dates, dimensions and evidence. An aggregation mapping
does not silently become a measurement conversion.

Standards differ in recognition, measurement, presentation and disclosure.
Consequently, conversion between IFRS, US GAAP and Chinese Accounting Standards
for Business Enterprises (ASBE) is not generally a dictionary lookup.
The original basis remains preserved, and any comparison basis is a separate
derived book. The same principle applies to tax and prudential capital. A tax
loss carryforward, accounting deferred tax asset and regulatory capital
deduction can concern the same economic loss while belonging to different
measurement systems.

\section{Time, revisions and information availability}

At least four times matter. A machine may enter service in November; its cost
appears in the accounts for the year ending December; the annual report is
published in March; and AccDSL ingests the report in April. A later report may
revise the December comparative. These times answer different questions and
must remain distinct.

The storage model uses validity intervals for economic assertions and
transaction intervals for database versions, following the distinction in
\citet{snodgrass1986}. Publication and availability timestamps are additional
domain fields. Database transaction time records what AccDSL contained; it
cannot substitute for what investors could have known. An ``as known on''
query applies an explicit availability rule and does not include a later
restatement simply because its economic period is earlier.

AccDSL supports three named views: the original filing view; the latest
restated view, with a stated information cutoff; and a harmonised research
view, with explicit transformations. No view changes the stored source facts.
Corrections have links to the replaced assertion and a reason. Where two
sources conflict without a clear supersession relationship, both survive with
a conflict record.

Instant values and interval flows have different types. A balance at 30 June
is not a six-month flow. Calendar intervals use half-open boundaries
$[t_0,t_1)$ internally; a filing period ending 30 June maps to the next boundary
at 1 July in its reporting calendar. Quarter-only earnings may be derived from
year-to-date values only if the periods, basis, perimeter, units and revision
vintage are compatible. The derivation retains both source observations and
their correlated rounding uncertainty. Fifty-two/fifty-three-week calendars,
stub periods and reporting-date changes require explicit calendar objects.

Accounting-policy changes, estimate changes and error corrections also have
different temporal effects. The applicable rule package decides whether a
change is prospective, requires restatement, or adjusts opening balances.
The engine never infers that effect from a generic ``revision'' flag.

\section{The company boundary changes}

Legal entity, reporting group and economic exposure are separate dimensions.
Acquisitions, disposals, changes in control, joint arrangements, special-purpose
vehicles and non-controlling interests can change the reporting perimeter.
A group's increase in property need not be capital expenditure; it may reflect
an acquisition, a currency translation or a revaluation. The bridge between
opening and closing property must identify these movements before a scenario
model can extrapolate investment.

The company graph therefore stores dated ownership, control and counterparty
relationships. Control is an accounting conclusion with supporting evidence,
not a rule that blindly consolidates every majority-owned entity. Group
reporting combines entity books with a separate adjustment and elimination
book. Intragroup receivables and payables require matching counterparties,
currencies and dates. Unmatched eliminations remain unresolved; a balancing
``plug'' is not a recovered transaction.

Acquired opening balances, goodwill, purchase-price allocation, divested
operations and non-controlling interests each have their own event types.
Analyses can preserve reported scope or define a constant perimeter, but the
latter is a derived modelling assumption, not the original record.

\section{Further problems that the representation must handle}

Several difficulties extend beyond differences in standards, industries and
company labels. Recognition and settlement occur on different dates: a sale,
invoice, cash receipt and impairment can belong to four periods. Net balances
hide gross flows: repeated borrowing and repayment can leave the same closing
debt. Netting can conceal offsetting exposures and liquidity needs. Different
contractual maturities and repricing dates can share the same current carrying
amount. Guarantees, undrawn credit facilities, litigation and covenants can
be economically important before an asset or liability is recognised.

Foreign currency introduces transaction currency, functional currency and
presentation currency, together with different translation conventions for
stocks and flows. Currency translation effects must not be mistaken for cash
receipts. Inventory and depreciation depend on vintages or lots; an aggregate
stock and average age may be insufficient under a shock. Tax bases can evolve
differently from carrying amounts. Deferred tax recognition requires assumptions
about recoverability, rather than multiplication by a tax rate alone.

Public information is selectively disclosed and rounded. Duplicate XBRL facts
may agree, overlap within rounding bounds or contradict one another. A dash
can mean zero, immaterial, unavailable or inapplicable. Tables can mix units,
span pages and embed signs in parentheses. Text and table references can be
wrong in the source itself. AccDSL retains these distinctions and gives an
analyst a local repair path without pretending that extraction confidence
settles economic interpretation.

Finally, strategic behaviour matters. A firm may refinance, cut investment,
raise equity, change prices, breach a covenant or fail. A bank may preserve a
capital ratio by shrinking assets; an insurer may change reinsurance or asset
allocation. Such responses cannot be deduced from the balance-sheet identity.
They require explicitly calibrated or assumed policies, and alternative
policies must remain separate when the available record does not distinguish
them.
